\subsection{IV, troisième partie : images, degrés et platitude relative}
\label{note:iv3-geometry}
Dans la preuve de 10.2, l'image directe par $Y_0\to X_0$ est un
faisceau sur $X_0$, et le fermé obtenu dans $X$ est $Y$.
En 10.4.11, les points dont le corps résiduel a un degré divisant $d$
sont les noyaux des homomorphismes vers $\mathbf F_{p^d}$. Un même noyau
peut correspondre à plusieurs plongements du corps résiduel : pour
$\mathbf F_{p^e}$, $e\mid d$, il y a $e$ plongements, permutés par
Frobenius. L'identification porte donc sur les noyaux, sans affirmer une
bijection avec les homomorphismes.

Les générateurs homogènes de 11.2.9 sont de degré zéro. En 11.3.9,
la platitude est relative à $f_\lambda$ ; une platitude sur la base
n'implique pas la platitude comme module sur $\mathcal O_{X_\lambda}$.
L'homomorphisme descendu est $h_\mu:\mathcal F_\mu\to\mathcal G_\mu$.
La démonstration de 11.3.10 se rapporte à cette proposition elle-même.
En 11.3.16, la propriété du morphisme composé $g\circ f$ est celle
utilisée pour les ouverts considérés. En 11.3.17, la finitude qui
permet le choix d'un étage est celle de $C$ comme $A$-algèbre ; la
finitude comme $B$-algèbre seule ne fournit pas ce choix sur $A$.

\subsection{IV, 11.4.2 : une obstruction explicite à la platitude}
\label{note:iv3-nonflat}
Dans l'exemple, $\mathfrak J K=0$ et
$\mathfrak J\otimes_A K\simeq K\ne0$.
L'inclusion $\mathfrak J\hookrightarrow A$, tensorisée par $K$, devient
donc l'application nulle $K\to K$. Elle n'est pas injective.
Un module plat préserve l'injectivité par tensorisation ; cette
application fournit directement l'obstruction requise.
Cet argument ne suppose pas qu'un module plat arbitraire soit libre.

\subsection{IV, troisième partie : calculs locaux et dimensions}
\label{note:iv3-modules}
Les deux produits de 11.7.5 ont pour valeurs $2ue'$ et $2ue''$ ; le
facteur deux provient de la somme des deux termes croisés.
Dans 11.9.10, l'idéal $\mathfrak J'_h$ est engendré par
$s_1^m,\ldots,s_h^m$ ; aucun générateur $s_0$ n'intervient.
En 11.9.11, c'est le quotient
$\mathcal F/\mathfrak m_x^n\mathcal F$ qui a le support annoncé.
Un sous-module $\mathfrak m_x^n\mathcal F$ ne possède pas, en général,
ce support : il peut rester non nul loin de $x$.

La condition de dimension plate au plus $i$ de 12.3.2 est
$\operatorname{Tor}^B_j(M,N)=0$ pour $j>i$ et tout $N$.
Pour $i=0$, exiger aussi l'annulation de $\operatorname{Tor}_0$
forcerait $M=0$, alors qu'un module libre non nul est plat.
En 13.3.5, les nombres comparés sont
$\dim(\mathcal O_{Y_j,y})$, et non les anneaux locaux eux-mêmes.
Le point considéré en 14.5.8 est $x'$.

Dans le diagramme de 14.5.11.1, la section $s:S\to Y$ d'un
morphisme $g:Y\to S$ induit $s':X\to X\times_S Y$ ; la projection
$p_1:X\times_S Y\to X$ vérifie $p_1s'=1_X$.
Ces domaines déterminent les quatre flèches horizontales.
En 15.4.3, le morphisme de la fibre a pour source $U_y$,
ouvert auquel l'argument a été restreint.

\subsection{IV, troisième partie : renvois}
\label{note:iv3-refs}
La proposition dont l'application est discutée en 11.9.8 est 11.9.7.
Le lemme utilisé à la fin de 11.9.16 est 11.9.12.
La remarque invoquée au début de 15.3.4 est 15.2.4(v).
\par
